% FILE: 01_research_question.tex
% Project: livestream
% Thesis Role: real-time-process-broadcast-and-livestream-standards

\section{Research Question}
\label{sec:livestream-rq}

\subsection{Problem Statement}
\label{sec:livestream-rq-problem}

Process mining research has historically operated on \emph{post-hoc} event logs:
a process executes, its trace is recorded, and analysts subsequently apply discovery
or conformance algorithms to the collected data. This regime treats process evidence
as a retrospective artifact rather than a live observable. The problem addressed by
the livestream project is whether a software factory---specifically the
process-intelligence research foundry---can be instrumented to emit a continuous,
real-time, OCEL 2.0-conformant event stream such that process mining algorithms can
be applied \emph{as the factory operates}, without waiting for batch collection.

The technical challenge is non-trivial. The factory's events originate from
heterogeneous sources: Git commit history, AGI subagent conversation replays stored
as JSONL files, browser-side Petri net simulation events, and cryptographic ledger
blocks. Each source has a different emission cadence, schema, and object model.
Unifying these into a single OCEL 2.0 property graph stream---with proper
event-to-object (E2O) and object-to-object (O2O) links---requires a coherent
broadcast architecture and a standards-compliant telemetry schema. The Telemetry
Auditor audit (v30.1.2, 2026-06-01) confirmed that this unification is not yet
achieved: the browser visualizer telemetry is 0\% OCEL 2.0 compliant and scores
2/5 on XES (IEEE 1849-2016) due to missing \texttt{lifecycle:transition} attributes
and non-UTC timestamps.

\subsection{Old-Regime Assumption Challenged}
\label{sec:livestream-rq-old-regime}

The old regime assumes that process evidence is collected offline and analyzed
offline. This assumption is embedded in the standard XES workflow: events are
written to a file, the file is imported into a tool such as pm4py, and conformance
checking proceeds as a batch computation. The livestream project challenges this
assumption by asserting that a factory operating under van der Aalst AGI
livestreaming standards must emit evidence \emph{continuously} and that any
gap between emission and analysis represents an interval during which non-conforming
execution could proceed undetected. The MAPE-K autonomic controller implemented in
\texttt{autonomic.js} formalizes this challenge: the Monitor-Analyze-Plan-Execute
loop runs on every tick, not on every batch, and the EWMA drift detector (alpha\,=\,0.15,
LCL\,=\,0.920) raises an alert the moment conformance fitness crosses the lower
control limit. Under the old regime, such a crossing would only be discovered after
the fact.

\subsection{Hypothesis}
\label{sec:livestream-rq-hypothesis}

\textbf{[AUTHOR THESIS]} The hypothesis of this project is that a software factory
can be made continuously self-witnessing by: (1) instrumenting each source of
process events with an OCEL 2.0-conformant emitter; (2) routing all event streams
through a unified broadcast layer; (3) subjecting the unified stream to real-time
conformance checking against a declared Petri net model; and (4) governing the
factory's own behavior through an autonomic feedback loop that responds to
conformance drift. When this architecture is fully instantiated, the gap between
process execution and process evidence collapses to the latency of the broadcast
layer---measured in the Stream Director audit (v30.1.2) at 6--15\,ms per frame,
well within the 200\,ms bound specified by the livestreaming standard.

\subsection{Success Criteria}
\label{sec:livestream-rq-success}

Success requires satisfying all of the following gate criteria simultaneously:
\begin{enumerate}
  \item All four broadcast modules emit OCEL 2.0-conformant events with proper
    \texttt{ocel:activity}, \texttt{ocel:timestamp} (UTC ISO 8601),
    \texttt{ocel:vmap}, and \texttt{ocel:omap} fields, and the visualizer telemetry
    achieves $\geq$\,4/5 on the XES compliance checklist.
  \item The browser visualizer executes without DOM binding errors; the Stream
    Director remediation (v30.1.2) must survive an independent re-audit.
  \item The EWMA drift detector correctly identifies injected non-conforming traces
    and triggers the autonomic controller to execute a remediation plan.
  \item The SHA-256 cryptographic ledger chain ($H_i = \text{SHA-256}(i \| t \|
    \text{case\_id} \| \text{payload} \| H_{i-1})$) is verified tamper-evident
    by the Ledger Custodian agent.
  \item Conformance metrics against the declared Petri net model reach
    Fitness\,$\geq$\,0.95 and Precision\,$\geq$\,0.95 on live stream data,
    corroborating the offline conformance audit result of Fitness\,=\,0.9880
    and Precision\,=\,1.0000.
\end{enumerate}

As of 2026-06-01, criteria 3 and 4 are partially supported by the autonomic
controller implementation and the ledger tamper test, but criteria 1 and 2 remain
formally open due to the Telemetry Auditor FAIL verdict.
